\section*{ВВЕДЕНИЕ}
\addcontentsline{toc}{section}{ВВЕДЕНИЕ}

Современные программные системы всё чаще сталкиваются с необходимостью анализа и обработки текстовых данных: от валидации пользовательского ввода и разбора логов до поиска по большим корпусам документов. Универсальным инструментом для решения подобных задач служат регулярные выражения -- формальный язык описания шаблонов, позволяющий компактно выразить сложные критерии поиска и замены текста. Регулярные выражения встроены в большинство современных языков программирования и широко применяются в текстовых редакторах, системах сборки, анализаторах логов и средствах обработки данных.

Несмотря на широкую распространённость готовых решений, разработка собственной библиотеки регулярных выражений представляет значительный теоретический и практический интерес: она требует глубокого понимания теории формальных языков, алгоритмов разбора и сопоставления строк, а также проектирования переносимых программных интерфейсов. Создание такой библиотеки на языке Си с экспортируемым C API обеспечивает её применимость в широком спектре программных систем без зависимости от среды выполнения или стороннего окружения.

Наряду с библиотекой практическую ценность представляет поисковое приложение с графическим интерфейсом, непосредственно демонстрирующее возможности разработанной библиотеки. Интерактивный текстовый редактор с поддержкой поиска, замены и рекурсивного обхода директорий является типичным примером программного обеспечения, в котором эффективный механизм регулярных выражений оказывает прямое влияние на удобство работы пользователя.

\emph{Цель настоящей работы} -- проектирование и реализация программно-информационной поисковой системы с графическим интерфейсом на основе библиотеки регулярных выражений собственной разработки. Для достижения поставленной цели необходимо решить \emph{следующие задачи:}
\begin{itemize}
	\item провести анализ предметной области и методов реализации движков регулярных выражений;
	\item разработать концептуальную модель библиотеки регулярных выражений;
	\item спроектировать программную систему библиотеки регулярных выражений на языке Си;
	\item спроектировать переносимый C API библиотеки для её интеграции в программы на Си-подобных языках;
	\item реализовать библиотеку регулярных выражений средствами языка Си;
	\item спроектировать программную систему поискового приложения с графическим интерфейсом;
	\item реализовать поисковое приложение на языке C++ с использованием фреймворка Qt.
\end{itemize}

\emph{Структура и объем работы.} Отчёт состоит из введения, 4 разделов основной части, заключения, списка использованных источников, 2 приложений. Текст выпускной квалификационный работы равен \formbytotal{lastpage}{страниц}{е}{ам}{ам}.

\emph{Во введении} сформулирована цель работы, поставлены задачи разработки, описана структура работы, приведено краткое содержание каждого из разделов.

\emph{В первом разделе} на стадии описания технической характеристики предметной области приводится анализ регулярных выражений, методов их реализации, языка Си и C API, а также фреймворка Qt.

\emph{Во втором разделе} на стадии технического задания приводятся требования к разрабатываемой библиотеке регулярных выражений, её C API и поисковому приложению с графическим интерфейсом.

\emph{В третьем разделе} на стадии технического проектирования представлены проектные решения для библиотеки регулярных выражений: внутренние структуры данных, формальная грамматика поддерживаемого синтаксиса, алгоритмы компиляции и сопоставления, описание экспортируемого API, а также архитектура поискового приложения и его модулей.

\emph{В четвёртом разделе} на стадии рабочего проекта приводится описание реализации разработанных программных систем, проводится тестирование библиотеки регулярных выражений и поискового приложения.

В заключении излагаются основные результаты работы, полученные в ходе разработки.

В приложении А представлен графический материал. В приложении Б представлены фрагменты исходного кода. 
